\section{Dataset Characterization}
\label{sec:dataset}

\subsection{Input corpus: scraping and deduplication}
\label{sec:dataset-scraping}

\paragraph{Scraping.}
Source PBTs come from the \texttt{Benchify/realpbt} dataset~\cite{benchify2024realpbt}, assembled by a custom Python scraper that operates on the GitHub dependency graph of the Hypothesis library.
The scraper takes a pre-collected list of $\approx$54k dependent repositories, iterates over them asynchronously with rate-limit-aware GitHub REST API backoff, shallow-clones each repository, and uses ripgrep to locate files containing the \texttt{@given} decorator.
For each such file, it parses the Python AST to extract all PBT functions---identified as function definitions whose decorator list includes a call to \texttt{given}---and performs a transitive static call-graph traversal to recover the minimal closure of locally-defined helpers called by each PBT.
Repository metadata (name, URL, license, stars, forks, commit hash) and extracted PBTs (source code, dependency names and their source) are persisted to PostgreSQL; repos whose commit hash was already processed are skipped, enabling incremental re-scraping.
Only repositories whose license matches a permissive allowlist (MIT, Apache-2.0, BSD variants, ISC, Unlicense, CC0, BlueOak-1.0.0) contribute PBTs to the released dataset.

\paragraph{Fork deduplication.}
The raw scrape yielded 54,345 PBTs, but GitHub's dependency graph surfaces many forks of the same upstream (e.g., hundreds of forks of \texttt{pytorch} each carrying an identical Hypothesis test suite).
We deduplicate by grouping repositories whose names share the same base component, selecting a canonical representative per group (highest star count, breaking ties by PBT count), and retaining PBTs from canonical repos plus any PBTs unique to forks (identified by SHA-256 of whitespace-normalized code).
After fork deduplication the corpus contains 11,039 PBTs across 333 repositories.

\subsection{Post-production pipeline}
\label{sec:dataset-postprod}

After the formalization agent produces Lean artifacts, a multi-stage post-production pipeline processes the raw run outputs into a final curated dataset.
The stages are: merge $\to$ validate $\to$ unify $\to$ dedup $\to$ enrich $\to$ grade $\to$ metrics.
Each stage is an independent CLI command (\texttt{uv run postprod <stage>}) and the full pipeline state is checkpointed so it can be resumed after interruption.

\paragraph{Merge.}
Multiple independent benchmark runs (identified by timestamped run directories) are merged into a single JSONL file.
Each sample directory contains \texttt{datapoint.json} (source PBT metadata), \texttt{qa.json} (agent outputs and metrics), and \texttt{Spec.lean}/\texttt{Impl.lean} (generated Lean artifacts).
When the same \texttt{sample\_id} appears in multiple runs, the merge step keeps the highest-quality copy, scored by: (1)~success flag, (2)~overall structural faithfulness, (3)~theorem count, (4)~presence of unit tests, (5)~implementation autoformalization success, (6)~specification signature success.
Git commit hash, model, run timestamp, and Lean toolchain version are recorded as provenance fields.

\paragraph{Validate.}
Because an agent can produce code that typechecks in isolation but fails under a full project build, we validate compiled samples by instantiating each \texttt{Spec.lean}/\texttt{Impl.lean} pair into a \texttt{lake build} project and running the Lean compiler.
Validation runs up to 400 samples by default (yielding $\pm$5\% accuracy at 95\% CI) or all samples with \texttt{--all}.
Only samples that pass \texttt{lake build} are retained in the dataset; failures are logged but dropped.

\paragraph{Unify.}
Our dataset was produced in two pipeline generations (\emph{feb03} and \emph{apr08}) with schema drift between them.
The unify step resolves these differences: it renames \texttt{radon} $\to$ \texttt{realpbt\_radon}, backfills the structured \texttt{dependencies} list from feb03's legacy parallel JSON-stringified arrays (filling \texttt{null} for fields not captured in that schema), drops 100\%-null dead fields (\texttt{lean\_metrics.tests\_structure}, \texttt{lean\_metrics.tests\_complexity}, \texttt{metrics\_metadata.tests\_available}), and reassigns fresh monotone \texttt{sample\_id}s 1..N with feb03 rows first.
Each row carries a \texttt{run} tag indicating its source generation.

\paragraph{Code-level deduplication.}
The same PBT may have been formalized in multiple runs (e.g., because sampling overlap exists between run index ranges, or because a PBT appeared in both feb03 and apr08).
This step groups samples by SHA-256 of \texttt{realpbt\_code} and keeps the single best representative per group (using the same quality scoring as the merge step), separately tracking cross-run vs.~within-run duplicates for audit.

\paragraph{Enrich.}
Every sample is annotated with:
\texttt{code\_hash} (SHA-256 of \texttt{realpbt\_code}),
\texttt{is\_canonical} (whether it is the best formalization of its PBT),
\texttt{formalization\_rank} (1-based rank within the group, 1 = canonical),
\texttt{num\_formalizations} (how many runs produced a formalization of this PBT), and
\texttt{pareto\_dominated} (whether a strictly better sample exists on all five structural faithfulness sub-metrics: parameter coverage, type correspondence, strategy coverage, assertion coverage, and dependency coverage).
Stale v1 grader fields with inconsistent coverage are also stripped here, and columns are reordered for clean HuggingFace presentation.

\paragraph{Difficulty grading.}
Each sample is scored for formalization difficulty by Claude Haiku, producing a binary easy/hard label (\texttt{difficulty\_binary}) with a calibrated confidence score and a brief reasoning trace.
The grader uses a structured prompt including the PBT source code, the generated Lean spec, and contextual information about the formalization strategy.
Grading is resumable: a checkpoint file accumulates completed grades so a killed run restarts where it left off.

\paragraph{Lean metrics.}
A final metrics pass runs radon on the Python source and a custom Lean AST analyzer on \texttt{Spec.lean}/\texttt{Impl.lean} to populate \texttt{lean\_metrics} (total lines, sorry count, axiom count, spec structure and complexity fields including Halstead volume and difficulty, theorem/def counts).

\subsection{Corpus statistics}
\label{sec:dataset-stats}

Table~\ref{tab:dataset-stats} summarizes the post-deduplication input corpus.
The star and fork distributions are heavily right-skewed: the median repository has 0~stars, reflecting that most real-world Hypothesis users work on small or private-to-GitHub projects.
The median PBT is 13 lines with a cyclomatic complexity of 2, consistent with the concise, assertion-focused style typical of Hypothesis tests.

\begin{table}[t]
  \centering
  \caption{Summary statistics for the post-deduplication RealPBT input corpus.}
  \label{tab:dataset-stats}
  \begin{tabular}{lr}
    \toprule
    \textbf{Metric} & \textbf{Value} \\
    \midrule
    PBTs (pre-dedup, HuggingFace)  & 54,345 \\
    PBTs (post-dedup, working set) & 11,039 \\
    Unique repositories            & 333    \\
    Median PBTs per repo           & 6      \\
    Mean PBTs per repo             & 33.2   \\
    Max PBTs in a single repo      & 956    \\
    Repos with only 1 PBT          & 61     \\
    \midrule
    Median repo stars              & 0      \\
    Mean repo stars                & 39     \\
    Repos with $\geq$100 stars    & 10     \\
    Repos with $\geq$1{,}000 stars& 1      \\
    \midrule
    Median PBT LOC                 & 13     \\
    Mean PBT LOC                   & 17.2   \\
    Median cyclomatic complexity   & 2.0    \\
    Mean cyclomatic complexity     & 3.3    \\
    \midrule
    MIT-licensed repos             & 148    \\
    Apache-2.0-licensed repos      & 66     \\
    BSD-3-Clause-licensed repos    & 25     \\
    \bottomrule
  \end{tabular}
\end{table}

Figures illustrating the distributions of PBTs per repo, stars, forks, and complexity are generated by \texttt{comms/graphs/src/graphs/realpbt.py} and are reproducible via \texttt{make realpbt} in \texttt{comms/graphs/}.

Table~\ref{tab:dataset-comparison} compares the RealPBT input corpus with the contemporaneous Hypothesis Corpus of DeVoe et al.~\cite{devoe2026hypothesis} (see also Section~\ref{sec:related}).

\begin{table}[t]
  \centering
  \setlength{\tabcolsep}{4pt}
  \caption{Comparison of the \FVSpec RealPBT corpus with the Hypothesis Corpus 2026~\cite{devoe2026hypothesis}. Both datasets deduplicate repositories but differ in method and granularity.}
  \label{tab:dataset-comparison}
  \begin{tabular}{lcc}
    \toprule
    & \textbf{Ours (RealPBT)} & \textbf{Hypothesis Corpus 2026} \\
    \midrule
    Collection method              & dependency graph      & API string search \\
    PBTs (before test-level dedup) & 54{,}345 & 28{,}928 \\
    PBTs (after test-level dedup)  & 11{,}039 & not performed \\
    Repositories                   & 333      & 1{,}529 \\
    Repo dedup method              & fork/name grouping & MinHash ($J{>}0.75$, 30\% files) \\
    Test-level dedup               & SHA-256 of PBT code & none \\
    License filter                 & permissive only & none \\
    \midrule
    Dependency source code         & \checkmark & \texttimes \\
    Radon code metrics             & \checkmark & \texttimes \\
    NL summary                     & \checkmark & \checkmark \\
    Abstract pattern / domain tag  & \texttimes & \checkmark \\
    \texttt{@settings} config     & \texttimes & \checkmark \\
    Runtime \& entropy data       & \texttimes & \checkmark \\
    Line coverage                  & \texttimes & \checkmark \\
    Execution traces (500 runs)    & \texttimes & \checkmark \\
    \midrule
    Storage format                 & JSONL & SQLite (3 DBs, $\sim$1.4\,GB) \\
    \bottomrule
  \end{tabular}
\end{table}